
\section{Introduction}

Bayesian optimization techniques form a successful 
approach for optimizing black-box functions \cite{Brochu:2009}. The goal
of these methods is to find the global maximizer of a nonlinear and generally non-convex function $f$
whose derivatives are unavailable. Furthermore, the evaluations of $f$ are usually corrupted by noise and
the process that queries $f$ can be computationally or economically very expensive. To address these challenges,
Bayesian optimization devotes additional effort to modeling the unknown function $f$ and its behavior.
These additional computations aim to minimize the number of evaluations that are needed to find the global optima.

Optimization problems are widespread in science and engineering and as a result
so are Bayesian approaches to this problem.  Bayesian optimization has
successfully been used in robotics to adjust the parameters of a robot's
controller to maximize gait speed and smoothness~\citep{Lizotte:2007} as well as
parameter tuning for computer graphics~\cite{Brochu:2007}. Another example
application in drug discovery is to find the chemical derivative of a particular
molecule that best treats a given disease~\cite{negoescu:2011}.  Finally,
Bayesian optimization can also be used to find optimal hyper-parameter values
for statistical~\cite{Wang:2013} and machine learning techniques
\citep{Snoek:2012}.

As described above, we are interested in finding the global maximizer $\xopt=\argmax_{\x\in\X} f(\x)$
of a function $f$ over some bounded domain, typically $\X\subset\R^d$. We 
assume that $f(\x)$ can only be evaluated via queries to a black-box that
provides noisy outputs of the form $y_i\sim \Normal(f(\x_i), \sigma^2)$. We
note, however, that our framework can be extended to other non-Gaussian
likelihoods. In this setting, we describe a sequential search algorithm that,
after $n$ iterations, proposes to evaluate $f$ at some location $\x_{n+1}$. To
make this decision the algorithm conditions on all previous observations
$\D_n=\{(\x_1,y_1),\dots,(\x_n,y_n)\}$. After $N$ iterations the algorithm 
makes a final recommendation $\xrec_N$ for the global maximizer of the latent function $f$.  

We take a Bayesian approach to the problem described above and use a probabilistic model for the
latent function $f$ to guide the search and to select $\xrec_N$.
In this work we use a zero-mean Gaussian process (GP) prior for $f$
\cite{Rasmussen:2006}. This prior is specified by a positive-definite kernel
function $k(\x,\x')$. Given any finite collection of points
$\{\x_1,\dots,\x_n\}$, the values of $f$ at these points are jointly zero-mean
Gaussian with covariance matrix $\vK_n$, where $[\vK_n]_{ij}=k(\x_i,\x_j)$. For
the Gaussian likelihood described above, the vector of concatenated observations $\vy_n$
is also jointly Gaussian with zero-mean. Therefore, at any
location $\x$, the latent function $f(\x)$ conditioned on past observations
$\D_n$ is then Gaussian with marginal mean $\mu_n(\x)$ and variance $v_n(\x)$
given by
\begin{align}
    \mu_n(\x) 
    &= \vk_n(\x)\T(\vK_n+\sigma^2\vI)^{-1}\vy_n\,, &
    v_n(\x)
    &= k(\x,\x) - \vk_n(\x)\T(\vK_n+\sigma^2\vI)^{-1}\vk_n(\x)\,,
    \label{eq:gpvar}
\end{align}
where $\vk_n(\x)$ is a vector of cross-covariance terms between $\x$ and
$\{\x_1,\dots,\x_n\}$.

Bayesian optimization techniques use the above predictive distribution $p(f(\x)|\D_n)$
to guide the search for the global maximizer $\xopt$. In particular, 
$p(f(\x)|\D_n)$ is used during the computation of an acquisition function $\alpha_n(\x)$
that is optimized at each iteration to determine the next evaluation location $\x_{n+1}$.
This process is shown in Algorithm~\ref{alg:bo}. Intuitively, the acquisition function
$\alpha_n(\x)$ should be high in areas where the maxima is most likely to lie given the current data.
However, $\alpha_n(\x)$ should also encourage exploration of the search space to
guarantee that the recommendation $\xrec_N$ is a global optimum of $f$, not just a global
optimum of the posterior mean.  Several acquisition functions have been proposed
in the literature. Some examples are the probability of
improvement~\citep{Kushner:1964}, the expected improvement~\citep{Mockus:1978}
or upper confidence bounds~\citep{Srinivas:2010}.  Alternatively, one can
combine several of these acquisition functions~\citep{Hoffman:2011}.

\begin{figure}
\vspace{-1.5em}
\setlength{\tabcolsep}{0pt}
\begin{tabular}{p{0.5\textwidth} p{0.5\textwidth}}
\vspace{0pt}
\begin{minipage}{0.48\textwidth}
\begin{algorithm}[H]
    \caption{Generic Bayesian optimization}
    \label{alg:bo}
    {\footnotesize
    \begin{algorithmic}[1]
        \REQUIRE a black-box with unknown mean $f$
        \FOR{$n=1,\dots,N$}
            \STATE select $\x_n = \argmax_{\x\in\X}\alpha_{n-1}(\x)$
            \STATE query the black-box at $\x_n$ to obtain $y_n$
            \STATE augment data $\D_n = \D_{n-1} \cup \{(\x_n, y_n)\}$
        \ENDFOR
        \RETURN $\xrec_N=\argmax_{\x\in\X} \mu_N(\x)$
        \vspace{3.2em}
    \end{algorithmic}}
    \end{algorithm}
\end{minipage}
&
\vspace{0pt}
\begin{minipage}{0.48\textwidth}
    \begin{algorithm}[H]
    \caption{PES acquisition function}
    \label{alg:pes}
    {\footnotesize
    \begin{algorithmic}[1]
        \REQUIRE a candidate $\x$; data $\D_n$
        \STATE \tikzmark{a}sample $M$ hyperparameter values $\{\vpsi\i\}$
        \FOR{$i=1,\dots,M$}
            \STATE sample $f\i\sim p(f|\D_n, \vphi, \vpsi\i)$
            \STATE set $\xopt\i\gets\argmax_{\x\in\X} f\i(\x)$
            \STATE compute $\vm_0\i$, $\vV_0\i$ and $\widetilde\vm\i$, $\widetilde\vv\i$
            \STATE compute $v_n\i(\x)$ and $v_n\i(\x|\xopt\i)$
        \ENDFOR
        \RETURN $\alpha_n(\x)$ as in (\ref{eq:marginalizedObjective})
    \end{algorithmic}}
    \end{algorithm}
\end{minipage}
\end{tabular}
\begin{tikzpicture}[remember picture, overlay]
    \node
    [draw=green!80!black, very thick, dotted, rectangle, anchor=north west,
     minimum width=5.8cm,
     minimum height=2.05cm]
    (box) at ($(a) + (-0.5, 0.30)$) {};
    \node[text=green!80!black, rotate=-90, anchor=south] at (box.east) {precomputed};
\end{tikzpicture}
\vspace{-1em}
\end{figure}

The acquisition functions described above are based on optimistic estimates of
the latent function $f$ which implicitly trade off between exploiting the
posterior mean and exploring based on the uncertainty.  We instead follow the
approach described in \cite{HennigSchuler:2012} and aim to maximize the expected
gain of information on the posterior distribution of the global maximizer
$\xopt$.  In Section~\ref{sec:method} we derive a rearrangement of this
acquisition function and a corresponding approximation that we call Predictive
Entropy Search (PES). PES is more accurate than the approximation used in
\cite{HennigSchuler:2012}.  In Section~\ref{sec:experiments} we empirically
evaluate this claim on both synthetic and real-world problems and show that this
leads to real gains in performance.

